% GENERATED FILE. Edit canonical lessons, then run npm run generate:books.
% canonical-source-hash: fnv1a64:eff327d760ad32a4
% canonical-lessons: TA-C361-tirmani, TA-C361-paricili, TA-C361-alavukol, TA-C361-poruttam, TA-C361-campantam

\chapter{Weighing It Up}
\label{ch:ta-weighing-it-up}
\hlchaptermodality{361}

\begin{chapteropening}
\textbf{By the end of this chapter:} \emph{I can say that I weigh a request and decide.}

Complete the last lesson of chapter 361: I can say that I weigh a request and decide.
\end{chapteropening}

\section[\texorpdfstring{\ta{தீர்மானி} (tīrmāṉi)}{தீர்மானி (tīrmāṉi)}]{\ta{தீர்மானி} (tīrmāṉi) — to determine; to resolve}

\label{lesson:TA-C361-tirmani}



\begin{quote}
Before the new one: say the Tamil for blame, then the Tamil for an excuse.
\end{quote}

\subsection*{You'll want to know: \ta{தீர்மானி}}
\textbf{\ta{தீர்மானி}} — \emph{tīrmāṉi} — \textquotedblleft{}to determine; to resolve\textquotedblright{}.

\begin{quote}
\ta{அடுத்த} \ta{ஆண்டு} \ta{வெளிநாட்டில்} \ta{படிக்கத்} \ta{தீர்மானித்தேன்}.
\end{quote}

\begin{quote}
\emph{aṭutta āṇṭu veḷināṭṭil paṭikkat tīrmāṉittēṉ.}
\end{quote}

\begin{quote}
I decided to study abroad next year.
\end{quote}

Past (I …): \ta{தீர்மானித்தேன்}, \emph{tīrmāṉittēṉ}.

With an infinitive: \ta{படிக்கத்} \ta{தீர்மானி}, to decide to study. The noun \ta{தீர்மானம்} is a resolution.

\begin{cousinweb}
\textbf{\ta{தீர்மானி}} (\emph{tīrmāṉi}) is native Dravidian: DEDR 3278 lists Tamil \emph{tīrmāṉam} and \emph{tīrmāṉi} in the family of \emph{tīr}, \textquotedblleft{}to end, be completed, be settled\textquotedblright{}.
\end{cousinweb}

\begin{grammarlens}[title={What you've built}]
One more word for giving reasons and opinions.
\end{grammarlens}

\subsection*{Your turn}
\begin{itemize}
\raggedright
  \item \emph{Say it:} \emph{tīrmāṉi}
  \item \emph{Say it:} \emph{tīrmāṉi}, once more
  \item \emph{Say it:} \emph{aṭutta āṇṭu veḷināṭṭil paṭikkat tīrmāṉittēṉ.} — I decided to study abroad next year.
  \item \emph{Recall:} say the Tamil for blame, then the Tamil for an excuse, then say \emph{tīrmāṉi} again
\end{itemize}

\begin{tcolorbox}[breakable,colback=teal!4,colframe=teal!35!black,title={Before you move on}]
What does \ta{தீர்மானி} mean? (\textquotedblleft{}To determine; to resolve\textquotedblright{}.) Say it once more.
\end{tcolorbox}
\section[\texorpdfstring{\ta{பரிசீலி} (paricīli)}{பரிசீலி (paricīli)}]{\ta{பரிசீலி} (paricīli) — to look into; to consider}

\label{lesson:TA-C361-paricili}



\begin{quote}
Before the new one: say the Tamil for an excuse, then the Tamil for to determine.
\end{quote}

\subsection*{You'll want to know: \ta{பரிசீலி}}
\textbf{\ta{பரிசீலி}} — \emph{paricīli} — \textquotedblleft{}to look into; to consider\textquotedblright{}.

\begin{quote}
\ta{அதிகாரிகள்} \ta{எங்கள்} \ta{புகாரைப்} \ta{பரிசீலித்தார்கள்}.
\end{quote}

\begin{quote}
\emph{atikārikaḷ eṅkaḷ pukāraip paricīlittārkaḷ.}
\end{quote}

\begin{quote}
The officials looked into our complaint.
\end{quote}

Past (I …): \ta{பரிசீலித்தேன்}, \emph{paricīlittēṉ}.

Formal, for requests and applications: \ta{விண்ணப்பத்தைப்} \ta{பரிசீலி}.

\begin{cousinweb}
\textbf{\ta{பரிசீலி}} (\emph{paricīli}) is a verb built on Sanskrit \emph{pariśīlana}, \textquotedblleft{}contact; application to, study\textquotedblright{} (Monier-Williams), with the loan-verb ending -\ta{இ}.
\end{cousinweb}

\begin{grammarlens}[title={What you've built}]
One more word for giving reasons and opinions.
\end{grammarlens}

\subsection*{Your turn}
\begin{itemize}
\raggedright
  \item \emph{Say it:} \emph{paricīli}
  \item \emph{Say it:} \emph{paricīli}, once more
  \item \emph{Say it:} \emph{atikārikaḷ eṅkaḷ pukāraip paricīlittārkaḷ.} — The officials looked into our complaint.
  \item \emph{Recall:} say the Tamil for an excuse, then the Tamil for to determine, then say \emph{paricīli} again
\end{itemize}

\begin{tcolorbox}[breakable,colback=teal!4,colframe=teal!35!black,title={Before you move on}]
What does \ta{பரிசீலி} mean? (\textquotedblleft{}To look into; to consider\textquotedblright{}.) Say it once more.
\end{tcolorbox}
\section[\texorpdfstring{\ta{அளவுகோல்} (aḷavukōl)}{அளவுகோல் (aḷavukōl)}]{\ta{அளவுகோல்} (aḷavukōl) — a yardstick; a criterion}

\label{lesson:TA-C361-alavukol}



\begin{quote}
Before the new one: say the Tamil for to determine, then the Tamil for to look into.
\end{quote}

\subsection*{You'll want to know: \ta{அளவுகோல்}}
\textbf{\ta{அளவுகோல்}} — \emph{aḷavukōl} — \textquotedblleft{}a yardstick; a criterion\textquotedblright{}.

\begin{quote}
\ta{வெற்றிக்குப்} \ta{பணம்} \ta{மட்டுமே} \ta{அளவுகோல்} \ta{அல்ல}.
\end{quote}

\begin{quote}
\emph{veṟṟikkup paṇam maṭṭumē aḷavukōl alla.}
\end{quote}

\begin{quote}
Money alone is not the measure of success.
\end{quote}

Literally a measuring rod: \ta{அளவு}, measure, + a word for a stick.

\begin{cousinweb}
\textbf{\ta{அளவுகோல்}} (\emph{aḷavukōl}) is native: \emph{aḷavu}, \textquotedblleft{}measure, extent\textquotedblright{} (DEDR 295), and \emph{kōl}, \textquotedblleft{}stick, staff\textquotedblright{} (DEDR 2237): a measuring rod, and so a criterion.
\end{cousinweb}

\begin{grammarlens}[title={What you've built}]
One more word for giving reasons and opinions.
\end{grammarlens}

\subsection*{Your turn}
\begin{itemize}
\raggedright
  \item \emph{Say it:} \emph{aḷavukōl}
  \item \emph{Say it:} \emph{aḷavukōl}, once more
  \item \emph{Say it:} \emph{veṟṟikkup paṇam maṭṭumē aḷavukōl alla.} — Money alone is not the measure of success.
  \item \emph{Recall:} say the Tamil for to determine, then the Tamil for to look into, then say \emph{aḷavukōl} again
\end{itemize}

\begin{tcolorbox}[breakable,colback=teal!4,colframe=teal!35!black,title={Before you move on}]
What does \ta{அளவுகோல்} mean? (\textquotedblleft{}A yardstick; a criterion\textquotedblright{}.) Say it once more.
\end{tcolorbox}
\section[\texorpdfstring{\ta{பொருத்தம்} (poruttam)}{பொருத்தம் (poruttam)}]{\ta{பொருத்தம்} (poruttam) — a good match; suitability}

\label{lesson:TA-C361-poruttam}



\begin{quote}
Before the new one: say the Tamil for to look into, then the Tamil for a yardstick.
\end{quote}

\subsection*{You'll want to know: \ta{பொருத்தம்}}
\textbf{\ta{பொருத்தம்}} — \emph{poruttam} — \textquotedblleft{}a good match; suitability\textquotedblright{}.

\begin{quote}
\ta{இந்தச்} \ta{சட்டைக்கும்} \ta{அந்தக்} \ta{கால்சட்டைக்கும்} \ta{நல்ல} \ta{பொருத்தம்}.
\end{quote}

\begin{quote}
\emph{intac caṭṭaikkum antak kālcaṭṭaikkum nalla poruttam.}
\end{quote}

\begin{quote}
This shirt and those trousers are a good match.
\end{quote}

\ta{பொருத்தமான}, suitable: \ta{பொருத்தமான} \ta{பதில்}, a fitting answer.

\begin{cousinweb}
\textbf{\ta{பொருத்தம்}} (\emph{poruttam}) is native Dravidian: DEDR 4541 lists Tamil \emph{poru}, \textquotedblleft{}to join, unite\textquotedblright{}, \emph{poruttu}, \textquotedblleft{}to fit together\textquotedblright{}, and \emph{poruttam}, \textquotedblleft{}fitness, suitability\textquotedblright{}.
\end{cousinweb}

\begin{grammarlens}[title={What you've built}]
One more word for giving reasons and opinions.
\end{grammarlens}

\subsection*{Your turn}
\begin{itemize}
\raggedright
  \item \emph{Say it:} \emph{poruttam}
  \item \emph{Say it:} \emph{poruttam}, once more
  \item \emph{Say it:} \emph{intac caṭṭaikkum antak kālcaṭṭaikkum nalla poruttam.} — This shirt and those trousers are a good match.
  \item \emph{Recall:} say the Tamil for to look into, then the Tamil for a yardstick, then say \emph{poruttam} again
\end{itemize}

\begin{tcolorbox}[breakable,colback=teal!4,colframe=teal!35!black,title={Before you move on}]
What does \ta{பொருத்தம்} mean? (\textquotedblleft{}A good match; suitability\textquotedblright{}.) Say it once more.
\end{tcolorbox}
\section[\texorpdfstring{\ta{சம்பந்தம்} (campantam)}{சம்பந்தம் (campantam)}]{\ta{சம்பந்தம்} (campantam) — a connection; relevance}

\label{lesson:TA-C361-campantam}



\begin{quote}
Before the new one: say the Tamil for a yardstick, then the Tamil for a good match.
\end{quote}

\subsection*{You'll want to know: \ta{சம்பந்தம்}}
\textbf{\ta{சம்பந்தம்}} — \emph{campantam} — \textquotedblleft{}a connection; relevance\textquotedblright{}.

\begin{quote}
\ta{எனக்கும்} \ta{அந்தச்} \ta{சம்பவத்துக்கும்} \ta{எந்தச்} \ta{சம்பந்தமும்} \ta{இல்லை}.
\end{quote}

\begin{quote}
\emph{eṉakkum antac campavattukkum entac campantamum illai.}
\end{quote}

\begin{quote}
I have nothing to do with that incident.
\end{quote}

X-\ta{க்கும்} Y-\ta{க்கும்} \ta{சம்பந்தம்}, a connection between X and Y.

\begin{cousinweb}
\textbf{\ta{சம்பந்தம்}} (\emph{campantam}) is borrowed from Sanskrit \emph{sambandha}, \textquotedblleft{}binding together, close connection\textquotedblright{} (Monier-Williams). Tamil writes the b and dh of the loan with \ta{ப} and \ta{த}.
\end{cousinweb}

\begin{grammarlens}[title={What you've built}]
One more word for giving reasons and opinions.
\end{grammarlens}

\subsection*{Your turn}
\begin{itemize}
\raggedright
  \item \emph{Say it:} \emph{campantam}
  \item \emph{Say it:} \emph{campantam}, once more
  \item \emph{Say it:} \emph{eṉakkum antac campavattukkum entac campantamum illai.} — I have nothing to do with that incident.
  \item \emph{Recall:} say the Tamil for a yardstick, then the Tamil for a good match, then say \emph{campantam} again
\end{itemize}

\begin{tcolorbox}[breakable,colback=teal!4,colframe=teal!35!black,title={Before you move on}]
What does \ta{சம்பந்தம்} mean? (\textquotedblleft{}A connection; relevance\textquotedblright{}.) Say it once more.
\end{tcolorbox}
